% CECI EST UN FICHIER DE MODÈLE LATEX POUR LES ARTICLES INCLUS DANS
% *Anthology of Computers and the Humanities*. AJOUTEZ L’OPTION
% 'final' LORS DE LA CRÉATION DE LA VERSION FINALE DE L’ARTICLE. 
% NE PAS modifier la classe du document
\documentclass[final,fra]{anthology-ch} % pour la version finale
%\documentclass[fra]{anthology-ch}         % pour la soumission

% CHARGER LES PACKAGES LaTeX
\usepackage{booktabs}
\usepackage{graphicx}

%Update

% Insertion de common.latex en enlevant babel
% https://github.com/jgm/pandoc-templates/blob/master/common.latex






%/ common.latex inclusion

% AJOUTEZ vos propres packages avec \usepackage{}

% TITRE DE LA SOUMISSION
% Remplacez ceci par le titre de votre soumission
\title{Éditions numériques : on cherche un responsable ! Qui pour les préserver une fois les projets terminés, et comment ?}

% INFORMATIONS SUR LES AUTEURS ET LES AFFILIATIONS
% Pour chaque auteur, utilisez une nouvelle commande \author, avec
% les numéros entre crochets indiquant les affiliations correspondantes
% (section suivante) et l’ORCID-ID pour chaque auteur.  

\author[1,3]{Edgar Lejeune}[
  orcid= https://orcid.org/0000-0001-7930-1866,
  email= edgar.lejeune@univ-rouen.fr
]


\author[1,2]{Marcello Vitali-Rosati}[
  orcid= https://orcid.org/0000-0001-6424-3229,
  email= marcello.vitali.rosati@umontreal.ca
]

\author[1]{Tony Gheeraert}[
  orcid= https://orcid.org/0000-0002-1302-8278,
  email= tony.gheeraert@univ-rouen.fr
]

\author[1,2]{Federico Siragusa}[
  orcid= https://orcid.org/0009-0002-8842-7043,
  email= federico.siragusa@umontreal.ca
]

\author[1]{Josselin Morvan}[
  orcid= ,
  email= josselin.morvan@univ-rouen.fr
]

\author[1,2]{Nolwenn Pamart}[
  orcid= https://orcid.org/0009-0005-9063-7404,
  email= nolwenn.pamart@univ-rouen.fr
]

\author[1,2]{Antoine Marchand}[
  orcid= ,
  email= antoinemarchand87@gmail.com
]

\author[1,2]{Clara Grometto}[
  orcid= https://orcid.org/0009-0001-2687-342X,
  email= clara.grometto@umontreal.ca
]

\affiliation{1}{Chaire d'Excellence en Édition Numérique, UR 3229 CEREdI, Université de Rouen-Normandie, France}
\affiliation{2}{Université de Montréal, Canada}
\affiliation{3}{UMR 7219 SPHERE, Université Paris-Cité, CNRS, France}


% MOTS-CLÉS
% Fournissez un ou plusieurs mots-clés séparés par des virgules
% en utilisant la commande suivante

\keywords[fra]{éditions numériques, durabilité, presses universitaires, droit du numérique}
\keywords[eng]{digital editions, durability, university press, digital law}

% MÉTADONNÉES DE LA PUBLICATION
% Ces champs seront remplis lors de la publication ; 
% ils peuvent être laissés par défaut lors de la soumission
\pubyear{2026}
\pubvolume{4}
\pagestart{99}
\pageend{102}
\conferencename{Actes de la Conférence Humanistica}
\conferenceeditors{Serena Crespi, Simon Gabay, Martin Grandjean, Ariane Pinche, Marie Puren et Léa Saint-Raymond}
\doi{10.63744/lWYDTk12UMmf}  
\customhead{}

\addbibresource{bibliography.bib}

%%%%%%%%%%%%%%%%%%%%%%%%%%%%%%%%%%%%%%%%%%%%%%%%%%%%%%%%%%%%%%%%%%%%%%%%%%%
% DÉBUT DU TEXTE
\paperorder{13}
\begin{document}

\maketitle

\begin{abstract}
How can we prevent system administrators from unpluging critical digital
editions? And who should be responsible for making such decisions? We
were confronted with the urgency of this question in October 2024, when
a critical digital edition of \emph{L'Astrée}--the result of work begun
in 2004--was taken down from the host university's servers. This
presentation outlines the various challenges we faced in saving this
``legacy'' digital edition, from negotiations with Tufts' IT department
to the retroconversion of the edition and its inclusion in the catalog
of the \emph{Presses Universitaires de Rouen et du Havre}. Based on this
example, we propose an organizational framework and pragmatic solutions
applicable to other contexts. We also use this specific case to address
broader issues about digital preservation and sustainability.
\end{abstract}

\section{Introduction}\label{introduction}

Comment empêcher les administrateurs système de «\,débrancher\,» les
éditions critiques numériques~? Et qui devrait être chargé de prendre de
telles décisions~? Nous avons été confrontés à l'urgence de cette
question en octobre 2024, lorsque le travail éditorial mené pendant
vingt ans par Eglal et Fékri Henein sur le roman français du \textsc{xvii}\textsuperscript{e}~siècle \emph{L'Astrée} a 
disparu des serveurs de l'université Tufts.
Cent quatre-vingt pages HTML en accès ouvert, issues d'un travail
critique inestimable sur l'un des romans les plus commentés de l'époque
moderne en Europe occidentale, ont failli disparaître en raison de
failles de sécurité dans les protocoles PHP. Fékri Henein et les
services informatiques de Tufts ont accepté de migrer le site vers les
serveurs de l'université de Rouen. Nous avons ensuite proposé, avec
l'équipe du CEEN, de convertir l'édition critique en XML-TEI et
d'assurer la préservation du travail d'Eglal et Fékri Henein avec l'aide
des Presses universitaires de Rouen et du Havre.

Cette communication présente les différents défis auxquels nous avons
été confrontés pour sauver cette édition numérique «\,ancienne\,»,
depuis les négociations avec les services informatiques de Tufts jusqu'à
la rétroconversion de l'édition et à la négociation de son intégration
dans le catalogue des Presses universitaires de Rouen et du Havre. Sur
la base de cet exemple, nous proposons un cadre organisationnel et des
solutions pragmatiques applicables à d'autres contextes. Nous mobilisons
également ce cas spécifique pour aborder des questions plus générales~:
comment s'engager dans la durabilité des anciennes éditions critiques
numériques~? Quelles sont les conditions \emph{sine qua non} pour
garantir qu'un projet «\,terminé\,» demeure accessible sur le long terme~? Et que nous apprend ce cas sur la durabilité des éditions numériques
en général, ainsi que sur le rôle des infrastructures séculaires telles
que les universités dans ce processus~?

\section{État de l'art}\label{uxe9tat-de-lart}

Depuis la fin des années 1990, la question de la durabilité des
productions en humanités numériques est devenue cruciale.

Les spécialistes ont souligné la volatilité des technologies
du Web et les difficultés majeures à rendre un projet durable à long
terme \autocite{cummings_academics_2023}. Une solution commode a
consisté à élaborer des lignes directrices et des principes permettant
de «~bien terminer~» un projet (voir
\href{https://endings.uvic.ca/outputs.html}{la bibliographie du projet
Endings}).

Toutefois, ces solutions supposent de s'engager dès le début du projet
d'édition numérique dans des principes techniques et organisationnels
précis, par exemple dans le choix d'un cadre ou d'une architecture
donnée (voir par exemple Smithies et al. \cite{smithies_managing_2019}). Elles
impliquent également l'adoption de pratiques spécifiques, telles que
l'usage de dépôts, par exemple Git, Zenodo ou Nakala. Le recours à des
infrastructures dédiées et excessivement généralistes a cependant été
critiqué ces dernières années
\cite{van_zundert_if_2012,viglianti_4_2025}.

\section{Corpus mobilisé}\label{corpus-mobilisuxe9}

Que faire alors lorsqu'il s'agit de préserver un projet déjà terminé~?
Qui doit être responsable de rendre une telle édition numérique
accessible sur le long terme~? Sous quel format (données seules~;
données et interface~; données, interface et documentation)~? Pour
quelle durée~? Et quelles sont les conditions techniques,
organisationnelles et juridiques permettant d'atteindre un tel
objectif~?

Le cas de l'édition d'Eglal et Fékri Henein est éclairant parce que~:
a)~il s'agissait d'une édition déjà achevée et publiée lorsque nous
avons dû en assumer la responsabilité~; b)~elle provenait d'une période
où aucune norme ni règle ne régissait clairement les pratiques d'édition
numérique~; et c)~aucune institution n'était considérée comme
responsable de sa durabilité au moment où nous avons engagé le projet.

Ce sont là les principales raisons pour lesquelles nous estimons que ce
cas soulève des questions fondamentales pour la durabilité des éditions
critiques numériques de manière générale. Une telle situation --- être
appelé à maintenir et à pérenniser un projet initié au début des années
2000 --- n'est, à notre connaissance, pas particulièrement courante pour
une équipe de recherche en humanités numériques aujourd'hui. Elle est
toutefois appelée à devenir de plus en plus fréquente.

\section{Méthodologie et description du protocole
suivi}\label{muxe9thodologie-et-description-du-protocole-suivi}

Le travail que nous avons lancé comporte cinq étapes : 1)~l'archivage de
la version existante de l'édition sur les serveurs de l'université de
Rouen~; 2)~la rétroconversion de l'édition HTML en XML-TEI, à l'aide de
feuilles de style de transformation XSLT et de corrections manuelles~;
3)~la conduite d'une évaluation appropriée de cette nouvelle version par
des spécialistes en informatique et des spécialistes de la littérature
française du \textsc{xvii}\textsuperscript{e}~siècle, en vue de son intégration dans le catalogue
des Presses universitaires de Rouen et du Havre~; 4)~l'attribution
d'ISSN et de DOI à l'édition afin de satisfaire aux exigences du dépôt
légal pour les publications électroniques~; enfin 5)~sa mise à
disposition dans les outils de recherche des bibliothèques
universitaires. Bien qu'elle s'aligne sur le processus habituel de
publication imprimée, cette approche a soulevé une série de questions
redoutables et étroitement imbriquées.

Cette communication défend le rôle central des presses universitaires
dans la pérennisation des éditions numériques. Les presses universitaires
peuvent s'appuyer sur des institutions séculaires, dotées
d'infrastructures solides (services informatiques, bibliothèques). Elles
ont également pour fonction de légitimer la production des savoirs, et
sont en mesure d'établir des contrats durables entre les universités et
les producteurs de ces éditions \autocite{epron_iii_2018}.

\section{Résultats et discussion}\label{ruxe9sultats-et-discussion}

Les résultats de notre travail seront présentés selon deux angles
différents.

Premièrement, dans quelle mesure produisons-nous une nouvelle édition en
rétroconvertissant l'édition hypertextuelle HTML en XML-TEI~? S'agit-il
simplement d'une forme de «\,réimpression\,» numérique~? Devons-nous
enrichir l'annotation sémantique~? Pouvons-nous transformer
l'architecture de l'édition et les nombreuses relations entre les
multiples pages HTML~? Les problèmes liés à la modélisation du texte
seront discutés en détail, d'une part, parce qu'ils sont décisifs pour
envisager la survie à long terme des éditions critiques numériques, et,
d'autre part, parce qu'ils révèlent des enjeux organisationnels et
juridiques clefs.

La deuxième série de questions auxquelles nous avons été confrontés
concerne la légitimation \emph{a posteriori} du travail éditorial mené
par Eglal et Fekri Henein. L'édition électronique de \emph{L'Astrée}
doit, à terme, être considérée comme un objet éditorial d'une valeur
équivalente à celle d'un livre imprimé. Les presses universitaires
seraient légalement tenues de la conserver et de la rendre accessible.
Pour ce faire, il convient de soumettre l'édition de \emph{L'Astrée} à
un protocole d'évaluation comparable à celui existant pour les éditions
imprimées. Ce protocole n'était, début 2025, pas encore mis en place aux
Presses universitaires de Rouen et du Havre. Nous nous sommes depuis
engagés à la fois dans la création d'un protocole d'évaluation dédié aux
éditions numériques, et dans l'élaboration d'une version de l'édition de
\emph{L'Astrée} susceptible d'être acceptée.

Nous avons constaté à cette occasion qu'aucun cadre juridique clair
n'existait aux Presses universitaires de Rouen et du Havre pour
s'engager auprès des auteurs de manière à garantir la durabilité d'une
édition critique numérique. Comment l'université peut-elle instituer
juridiquement l'obligation d'assurer la perennité d'une édition, à la
fois du point de vue de sa conservation et de son accessibilité~? Cette
obligation concerne-t-elle uniquement les données, ou également
l'interface du site web~? Et quels sont les moyens matériels, les
équipements, et les ressources humaines nécessaires pour que les presses
universitaires soient en mesure de remplir leurs obligations~? Cette
communication proposera des réponses à ces questions sur la base des
résultats d'une enquête préliminaire menée en France auprès de plusieurs presses universitaires, ainsi que de la littérature existante sur le droit des éditions numériques.

\section{Perspectives}\label{perspectives}

Comme nous le montrerons, l'expérience que nous avons menée à Rouen au
contact de l'édition critique numérique de \emph{L'Astrée} soulève un
ensemble de questions plus générales pour la durabilité et
l'accessibilité des éditions numériques. Bien que limitée au contexte
spécifique de la France, elle ouvre d'importantes discussions sur la
fonction --\,et la nécessaire transformation\,-- des infrastructures
séculaires telles que les presses universitaires. Les solutions que nous
avons mises en œuvre susciteront, nous l'espérons, l'intérêt de nos
collègues d'autres universités, dans d'autres contextes. Nous sommes
convaincus qu'elles donneront lieu à des débats féconds, et à
l'ouverture de collaborations nouvelles.

\printbibliography



% Vous pouvez inclure une annexe en utilisant la commande suivante
%\appendix

%\section{Première section de l’annexe} \label{appdx:first}


\end{document}